% Copyright (C) 2026 Christopher Hoover (ch@murgatroid.com)
% SPDX-License-Identifier: CC-BY-NC-SA-4.0
% Licensed under Creative Commons Attribution-NonCommercial-ShareAlike 4.0 International
% https://creativecommons.org/licenses/by-nc-sa/4.0/
\documentclass[12pt]{article}
\usepackage[T1]{fontenc}
\usepackage{lmodern}
\usepackage{amsmath}
\usepackage{amssymb}
\usepackage{booktabs}
\usepackage{tabularx}
\usepackage{float}
\usepackage{graphicx}
\usepackage{geometry}
\usepackage{array}
\usepackage{siunitx}
\usepackage{microtype}
\usepackage[section]{placeins}
\usepackage[hidelinks]{hyperref}
\usepackage{bookmark}

\geometry{margin=1in}

\hypersetup{
    pdftitle={A Diode Is Not a Square-Law Device},
    pdfauthor={Christopher Hoover (AI6KG, ch@murgatroid.com)},
    pdfsubject={Diode square-law approximation, detection, mixing}
}

% Upright roman subscripts (IEEE / IUPAP convention)
\newcommand{\ID}{I_{\mathrm{D}}}
\newcommand{\VD}{V_{\mathrm{D}}}
\newcommand{\IS}{I_{\mathrm{S}}}
\newcommand{\VT}{V_{\mathrm{T}}}
\newcommand{\VQ}{V_{\mathrm{Q}}}
\newcommand{\IQ}{I_{\mathrm{Q}}}
\newcommand{\vd}{v_{\mathrm{d}}}
\newcommand{\id}{i_{\mathrm{d}}}
\newcommand{\gd}{g_{\mathrm{d}}}
\newcommand{\rd}{r_{\mathrm{d}}}
\newcommand{\vhd}{\hat{v}_{\mathrm{d}}}
\newcommand{\dd}{\mathrm{d}}
\newcommand{\eu}{\mathrm{e}}

\renewcommand{\topfraction}{0.85}
\renewcommand{\bottomfraction}{0.7}
\renewcommand{\textfraction}{0.1}
\renewcommand{\floatpagefraction}{0.7}

\title{A Diode Is Not a Square-Law Device}
\author{Christopher Hoover, AI6KG \\ \texttt{ch@murgatroid.com}\thanks{\copyright{} 2026 Christopher Hoover. Licensed under CC BY-NC-SA 4.0: \url{https://creativecommons.org/licenses/by-nc-sa/4.0/}}}
\date{}

\begin{document}

\maketitle

\begin{abstract}
The diode is routinely described as a ``square-law device''---a useful shorthand for explaining RF detection and frequency mixing. Its I--V characteristic is exponential, not quadratic, and the square-law approximation holds only over roughly one decade of signal amplitude. This article derives the four operating regimes that emerge from a Taylor expansion of the Shockley equation, establishes quantitative amplitude boundaries between them using second- and third-harmonic distortion criteria, and traces the practical consequences for detector sensitivity, mixer design, and LO drive level.
\end{abstract}

Pick up almost any RF textbook and you will find the diode described as a ``square-law device''---invoked to explain crystal radio detection, mixer frequency conversion, and intermodulation analysis alike.

There is just one problem: it is not true.

A diode's current--voltage characteristic is exponential, not quadratic.

\section*{The Shockley equation}

The current through an ideal p--n junction diode is given by
\begin{equation}
\ID = \IS \left( \eu^{\VD / n \VT} - 1 \right),
\label{eq:shockley}
\end{equation}
where $\IS$ is the reverse saturation current, $n$ is the ideality factor (a device-dependent parameter, typically 1--2), $\VD$ is the junction voltage, and $\VT = kT/q \approx \SI{25.85}{\milli\volt}$ at $T = \SI{300}{\kelvin}$ is the thermal voltage.

For forward bias well above the thermal voltage---specifically, $\VD \gtrsim 4n\VT$, where the $-1$ term contributes less than \SI{2}{\percent} to $\ID$---the exponential approximation holds:
\begin{equation}
\ID \approx \IS\, \eu^{\VD / n \VT}.
\label{eq:exp-approx}
\end{equation}
On a linear scale the I--V curve shoots nearly straight up once conduction begins; on a semilog plot it is a straight line spanning many decades of current. It is not a parabola.

\section*{The power-series expansion}

To understand small-signal behavior, decompose the applied voltage into a DC bias $\VQ$ and a time-varying perturbation $\vd(t)$:
\[
  \VD = \VQ + \vd.
\]
Substituting into \eqref{eq:exp-approx} and defining the Q-point current $\IQ \equiv \IS\, \eu^{\VQ / n \VT}$,
\[
  \ID = \IQ\, \eu^{\vd / n \VT}.
\]
Expanding in a Maclaurin series and collecting the AC component $\id \equiv \ID - \IQ$, with small-signal conductance $\gd \equiv \IQ / (n\VT) = 1/\rd$:
\begin{equation}
  \id = \gd \vd + \frac{\gd}{2\, n \VT}\, \vd^{\,2} + \frac{\gd}{6\,(n \VT)^2}\, \vd^{\,3} + \cdots.
  \label{eq:power-series}
\end{equation}
The first term is linear (resistive), the second quadratic, the third cubic. The series converges for all $\vd$; how many terms are needed depends on the signal amplitude relative to $n\VT$.

\subsection*{Harmonic distortion as a yardstick}

For a sinusoidal drive $\vd = \vhd \cos\omega t$, each nonlinear term in \eqref{eq:power-series} generates harmonics of $\omega$. The second- and third-harmonic distortions relative to the fundamental are
\begin{equation}
  \mathrm{HD2} \approx \frac{\vhd}{4\, n \VT}, \qquad
  \mathrm{HD3} \approx \frac{\vhd^{\,2}}{24\,(n \VT)^2}.
  \label{eq:hd}
\end{equation}
These ratios define quantitative boundaries between operating regimes. HD2 reaches \SI{1}{\percent} at $\vhd \approx 0.04\,n\VT$ and \SI{25}{\percent} at $\vhd \approx n\VT$; HD3 reaches \SI{4}{\percent} at $\vhd = n\VT$. The \SI{1}{\percent} and \SI{25}{\percent} thresholds are useful engineering markers---not physical phase transitions---but they define a well-motivated one-decade window within which the two-term approximation is a fair description.

\section*{Why ``square law'' seems right}

Within the regime $0.04\,n\VT \lesssim \vhd \lesssim n\VT$, only the first two terms of \eqref{eq:power-series} matter:
\begin{equation}
  \id \approx \gd \vd + \frac{\gd}{2 n \VT} \vd^{\,2}.
  \label{eq:quadratic}
\end{equation}
The quadratic term is genuinely useful---but for different reasons depending on whether one or two tones are present.

\textbf{Detection (single tone).}
When an RF signal $\vd = \vhd \cos\omega t$ drives the diode, the quadratic term produces a DC component proportional to $\vhd^{\,2}$ and hence to signal power.
This is the operating principle of small-signal envelope detectors, crystal radio receivers, and microwave power detectors: the diode converts RF power directly to a DC current, and the output tracks power linearly---the defining property of a square-law detector.

\textbf{Mixing (two tones).}
When two signals at $\omega_1$ and $\omega_2$ are simultaneously present, the quadratic term cross-multiplies them via $\cos\omega_1 t\cdot\cos\omega_2 t = \tfrac{1}{2}[\cos(\omega_1+\omega_2)t + \cos(\omega_1-\omega_2)t]$, producing outputs at the sum and difference frequencies $\omega_1 \pm \omega_2$.
The linear term passes both signals through unchanged and produces no conversion.
This is the basis of resistive RF mixers: the $v^2$ nonlinearity provides frequency translation.

\section*{The catch: the linear piece always wins}

The quadratic term is not the whole story. The linear term $\gd \vd$ is always present, and for small signals it always dominates.

The ratio of the quadratic term to the linear term is $\vd/(2n\VT)$. For signals well below $2n\VT \approx \SI{50}{\milli\volt}$ (for $n \approx 1$), the linear piece is larger. The diode is always more linear than it is square-law. Equivalently: the I--V curve near any DC bias point approximates a straight line---a slope, or conductance---and the quadratic term is the \emph{curvature} of that line. Curvature is always a small correction to slope for small perturbations.

Below the square-law range ($\vhd \lesssim 0.04\,n\VT \approx \SI{1}{\milli\volt}$ for $n \approx 1$), the diode is effectively a pure resistor $\rd = 1/\gd$: no rectification, no frequency conversion, no mixing. Detection sensitivity collapses and a mixer produces no IF output.

\section*{The four operating regimes}

As drive amplitude increases, the diode traverses four regimes in turn. Table~\ref{tab:diode-regimes} summarizes them; Figure~\ref{fig:regimes} maps them onto a 1N4148 at room temperature.

\subsection*{Linear ($\vhd \lesssim 0.04\,n\VT$)}

Only the first term of \eqref{eq:power-series} is significant: $\id \approx \gd \vd$.
The diode is a small-signal resistor of dynamic resistance $\rd = n\VT/\IQ$.
No mixing, no rectification, no harmonic generation occurs.
This is the regime of biasing calculations and Q-point dynamic-resistance measurements.

\subsection*{Quadratic / square-law ($0.04\,n\VT \lesssim \vhd \lesssim n\VT$)}

HD2 grows from \SI{1}{\percent} to \SI{25}{\percent}; HD3 reaches \SI{4}{\percent} at the upper edge but is still small enough to ignore.
The two-term approximation \eqref{eq:quadratic} is valid.
Single-tone detection gives DC output proportional to power; two-tone mixing produces $\omega_1 \pm \omega_2$ products with intermodulation dominated by third-order terms.

For a Ge or Schottky diode ($n \approx 1$) biased at $\IQ \approx \SI{100}{\micro\ampere}$, this range spans roughly \SI{1}{\milli\volt} to \SI{26}{\milli\volt} of junction amplitude---about one decade.

\subsection*{Intermediate exponential ($n\VT < \vhd \lesssim \VQ - 4n\VT$)}

Once $\vhd$ exceeds $n\VT$, the two-term series is no longer adequate, but the forward-exponential approximation \eqref{eq:exp-approx} still holds throughout the voltage cycle---the negative peak has not yet driven the junction below the $4n\VT$ validity threshold. Setting $x \equiv \vhd/(n\VT)$ and substituting $\vd = \vhd\cos\omega t$ into \eqref{eq:exp-approx}, the modified Jacobi--Anger identity gives an exact Fourier expansion:
\begin{equation}
  \ID(t) \approx \IQ \left[\, I_0(x) + 2 \sum_{k=1}^{\infty} I_k(x)\, \cos(k \omega t) \,\right],
  \label{eq:bessel}
\end{equation}
where $I_k$ are modified Bessel functions of the first kind \cite{companion}.

Two features distinguish this regime.
First, a single tone generates a rich harmonic spectrum \emph{without any clipping}---exponential asymmetry alone is sufficient.
Second, the average (DC) current increases with drive amplitude: $\langle\ID\rangle = \IQ\, I_0(x) > \IQ$ for $x > 0$.
This rectified DC component---and the bias shift it produces in a self-biased circuit---is the operating principle of large-signal envelope detectors and grid-leak / cathode-bias stages.

The intermediate regime exists only when $\VQ > 5n\VT$, so a Si p--n diode has a wide unclipped exponential range while a Ge point-contact diode with lower $\VQ$ has a narrow one (see Table~\ref{tab:diode-numbers}).

\subsection*{Full-Shockley / rectifying ($\vhd \gtrsim \VQ - 4n\VT$)}

When the negative half-cycle drives $\VD$ below $4n\VT$, the $-1$ in \eqref{eq:shockley} is no longer negligible and the full Shockley equation is required.
A reverse-biased interval appears when $\vhd \gtrsim \VQ$; the waveform approaches half-wave rectification when $\vhd \gtrsim \VQ + 4n\VT$.
This is the regime of large-signal AM envelope detection, power rectification, switching, and clamping.
Circuit effects---load impedance, charge-storage time constants, series resistance---typically dominate the waveform shape in real rectifier circuits.

\begin{figure}[H]
\centering
\includegraphics[width=0.95\linewidth]{regimes_figure.pdf}
\caption{Four operating regimes for a Si p--n signal diode (1N4148, $n=1.85$, bias current $\IQ \approx \SI{100}{\micro\ampere}$). (a)~Semilog $|\ID|$--$\VD$ curve from the full Shockley equation~\eqref{eq:shockley}; the unshaded region ($\VD \gtrsim 4n\VT$) is where the exponential approximation~\eqref{eq:exp-approx} holds. (b)~Regime boundaries $0.04\,n\VT$, $n\VT$, and $\VQ-4n\VT$ on a logarithmic signal-amplitude axis, spanning more than two decades.}
\label{fig:regimes}
\end{figure}

\FloatBarrier

\begin{table}[H]
\centering
\renewcommand{\arraystretch}{1.15}
\caption{Diode operating regimes expressed as algebraic functions of $n$, $\VT$, and $\VQ$. Boundaries for the first two rows are based on HD2 $=$ \SI{1}{\percent} and \SI{25}{\percent}; the last two are based on validity of the exponential approximation~\eqref{eq:exp-approx}.}
\label{tab:diode-regimes}

\vspace{0.5\baselineskip}

\newcolumntype{L}[1]{>{\raggedright\arraybackslash\hsize=#1\hsize}X}
\begin{tabularx}{\textwidth}{L{0.85} L{0.95} L{1.10} L{1.10}}
\toprule
\textbf{Regime} & \textbf{Amplitude criterion} & \textbf{Dominant behavior} & \textbf{Typical application} \\
\midrule
Linear &
$\vhd \lesssim 0.04\, n \VT$ \newline (HD2 $<$ \SI{1}{\percent}) &
$\id \approx \gd \vd$ &
Biasing analysis, dynamic-resistance calculations. \\

Quadratic / square-law &
$0.04\, n \VT \lesssim \vhd \lesssim n \VT$ \newline (HD2 from \SI{1}{\percent} to \SI{25}{\percent}) &
$\id \approx \gd \vd + \dfrac{\gd}{2 n \VT} \vd^{\,2}$ &
RF mixers, microwave power detectors, intermodulation analysis. \\

Intermediate exponential &
$n \VT < \vhd \lesssim \VQ - 4n\VT$ \newline (exp.\ approx.\ holds; truncation inadequate) &
Bessel expansion \eqref{eq:bessel}; rich harmonics without clipping; $\langle\ID\rangle = \IQ I_0(x)$ &
Large-signal envelope detectors, harmonic generators, mixer LO drive. \\

Full-Shockley / rectifying &
$\vhd \gtrsim \VQ - 4n\VT$ \newline (reverse interval for $\vhd \gtrsim \VQ$) &
Full Shockley \eqref{eq:shockley} required; asymmetric conduction &
AM detection, power rectification, switching, clamping. \\
\bottomrule
\end{tabularx}
\end{table}

\FloatBarrier

\section*{Numerical values for common small-signal diodes}

The boundaries in Table~\ref{tab:diode-regimes} are algebraic. Table~\ref{tab:diode-numbers} instantiates them for three common technologies at $T = \SI{300}{\kelvin}$ and $\IQ \approx \SI{100}{\micro\ampere}$.

\begin{table}[H]
\centering
\renewcommand{\arraystretch}{1.3}
\setlength{\tabcolsep}{5pt}
\caption{Square-law and intermediate-regime boundaries for representative small-signal diodes at $T = \SI{300}{\kelvin}$ and $\IQ \approx \SI{100}{\micro\ampere}$. The linear-regime upper edge ($0.04\,n\VT$) ranges from \SI{1.1}{\milli\volt} to \SI{1.9}{\milli\volt} across all three. Values are approximate, extracted from datasheet typical-curve plots and SPICE model parameters.}
\label{tab:diode-numbers}

\vspace{0.5\baselineskip}

\begin{tabular}{l c c c c}
\toprule
\textbf{Diode} & $n$ & $\VQ$ & $n\VT$ (sq-law upper) & $\VQ - 4n\VT$ (interm.\ upper) \\
\midrule
Si p--n signal (1N4148)        & 1.85 & \SI{0.55}{\volt} & \SI{48}{\milli\volt} & \SI{359}{\milli\volt} \\
Si Schottky (1N5711, BAT54)    & 1.08 & \SI{0.30}{\volt} & \SI{28}{\milli\volt} & \SI{188}{\milli\volt} \\
Ge point-contact (1N34A)       & 1.05 & \SI{0.20}{\volt} & \SI{27}{\milli\volt} & \SI{91}{\milli\volt}  \\
\bottomrule
\end{tabular}
\end{table}

Two observations stand out. First, the Si p--n device has roughly $1.7\times$ the square-law amplitude headroom of Ge or Schottky, because its higher ideality factor raises $n\VT$; a 1N4148 can tolerate a junction swing up to $\sim$\SI{48}{\milli\volt} before HD2 exceeds \SI{25}{\percent}, versus $\sim$\SI{27}{\milli\volt} for a 1N34A. Second, the lower $\VQ$ of Ge and Schottky means the intermediate exponential range is narrow: the upper edge sits near \SI{91}{\milli\volt} for Ge (about \SI{64}{\milli\volt} beyond the square-law boundary), versus \SI{359}{\milli\volt} for a Si p--n diode (about \SI{311}{\milli\volt} of unclipped exponential operation). This explains the historical division of labor: Ge and zero-bias Schottky for sensitive weak-signal detection where low $\VQ$ matters most, and Si p--n for switching and mixer service where the wider amplitude range is useful.

\section*{Practical consequences}

\textbf{Why Ge and Schottky beat Si for weak-signal detection.}
The square-law regime begins at $\vhd \approx 0.04\,n\VT$---a lower threshold in absolute terms for Ge and Schottky because their forward voltage is lower.
Less input power is needed to reach the regime where rectification is efficient.
This is the fundamental reason crystal radio builders and microwave detector designers reach for Ge and zero-bias Schottky diodes.

\textbf{Why bias helps.}
Biasing the diode at a small forward current $\IQ$ raises the conductance $\gd = \IQ/(n\VT)$ and hence the responsivity $\gd/(2n\VT)$ of the quadratic term.
The square-law range shifts to lower signal amplitudes, improving sensitivity at the cost of a bias circuit and quiescent current drain.

\textbf{LO drive level in mixers.}
A mixer driven within the square-law range converts signal power to IF power with predictable gain and third-order intermodulation products that follow the familiar two-tone IP3 model.
Driving too hard moves the LO junction into the intermediate exponential regime: conversion gain compresses and additional higher-order products appear.
``Square-law mixer'' behavior is valid only over a moderate range of LO junction amplitude.

\textbf{High-level switching mixers.}
A different class of mixers takes the opposite approach: the LO drives the diodes deep into the rectifying regime, reversing the signal path at the LO rate rather than exploiting a $v^2$ nonlinearity.
Diode ring and double-balanced mixers work this way.
They are \textbf{not square-law devices at all}; their conversion performance derives from the switching action and the symmetry of the topology.
They offer better dynamic range and conversion efficiency than square-law designs at the cost of higher LO drive power.

\textbf{Zero-bias detectors.}
Near $\VQ = 0$---the operating point of zero-bias Schottky detectors---the $-1$ term in \eqref{eq:shockley} is not negligible at the bias point and the regime boundaries above do not apply; a separate analysis expanding the full Shockley equation about $\VD = 0$ is required \cite{maas}.

\section*{Temperature effects}

Both $\VT$ and $\VQ$ shift with temperature, moving the regime boundaries.
$\VT$ scales linearly with absolute temperature ($\sim$\SI{20}{\milli\volt} at \SI{-40}{\degreeCelsius} to \SI{31}{\milli\volt} at \SI{85}{\degreeCelsius}), stretching the linear and square-law cutoffs by roughly $\pm\SI{20}{\percent}$; meanwhile $\VQ$ at constant bias falls at about \SI{-2}{\milli\volt\per\kelvin} for Si, compressing the intermediate-regime upper edge $\VQ - 4n\VT$ faster still since $4n\VT$ rises simultaneously.
A fixed-bias detector reads this drift directly as a DC baseline shift at its output.

\section*{The bottom line}

A diode is a square-law device in the same way a stretch of road is flat: it is a useful approximation over a limited range, and knowing the limits prevents you from driving off the edge.
The exponential is the truth; the square law is a convenience that holds over roughly one decade of junction amplitude---between \SI{1}{\percent} and \SI{25}{\percent} second-harmonic distortion.
Understanding that range, and what sits on either side of it, is what separates a detector that works from one that merely seems to.

\begin{thebibliography}{9}
\setlength{\itemsep}{2pt}
\setlength{\parsep}{0pt}

\bibitem{sze}
S.\ M.\ Sze and K.\ K.\ Ng, \textit{Physics of Semiconductor Devices}, 3rd ed.\ Hoboken, NJ: Wiley-Interscience, 2007.

\bibitem{maas}
S.\ A.\ Maas, \textit{Microwave Mixers}, 2nd ed.\ Norwood, MA: Artech House, 1993.

\bibitem{nxp1n4148}
Nexperia, \textit{1N4148; 1N4148W; 1N4448 --- High-speed diodes}, data sheet (Si p--n; $n \approx 1.85$, $\VQ \approx \SI{0.55}{\volt}$).

\bibitem{skyworks1n5711}
Skyworks Solutions, \textit{1N5711 --- Schottky Barrier Diode}, and Nexperia, \textit{BAT54 series}, data sheets (Si Schottky; $n \approx 1.08$, $\VQ \approx \SI{0.30}{\volt}$; BAT54 is the SMT equivalent of 1N5711).

\bibitem{1n34a}
NTE Electronics, \textit{NTE109 / 1N34A --- Germanium Signal Detector Diode}, data sheet (Ge point-contact; $n \approx 1.05$, $\VQ \approx \SI{0.20}{\volt}$).

\bibitem{companion}
C.\ Hoover, \textit{Derivation of Diode Small-Signal and Square-Law Behavior}, technical note, 2026. [Online]. Available: \url{https://charlieh0tel.github.io/ee_stuff/math/diode_small_signal/diode_smallsignal.pdf}

\end{thebibliography}

\end{document}
